\section{Introduction}

Firm size is a well-established correlate of wages in labor economics \citep{BrownMedoff1989,OiIdson1999,troske1999evidence}. In labor markets where many workers are own-account workers or employed in micro-firms, wage differences between small and large firms help identify where questions about firm growth, worker allocation, and wage inequality are most relevant. This paper asks whether Colombian workers in larger firms earn higher real hourly wages than observably similar workers in smaller firms, how much of that gap is explained by observable worker characteristics, whether the gap has increased or decreased over time, and whether the wage gap between small and large firms differs between formal and informal employment.

Our motivation to measure the firm-size wage premium in Colombia is both empirical and policy-oriented. A recent policy diagnosis argues that low-income households are disproportionately linked to own-account work, very small businesses, informality, and low-earning jobs, whereas higher-income households are more often connected to larger and more formal employers \citep{eslava2026crecimiento}. This diagnosis is suggestive, but it does not show whether firm size remains associated with higher wages after accounting for worker characteristics and formality. The observed link between larger employers and higher household income may instead reflect the fact that larger firms employ workers with different characteristics, or that larger firms are more likely to offer formal jobs. This distinction matters because the wage gap associated with firm size has not been systematically documented in Colombia. The question is especially relevant in Colombia because employment remains concentrated outside the largest-firm categories, while recent growth theory links the rise of large firms to economic development \citep{Chen2026}.

We study the firm-size wage premium in Colombia using harmonized microdata from the Gran Encuesta Integrada de Hogares (GEIH) for 2008--2019 and 2021--2025. We construct a comparable firm-size classification, measure real hourly labor income, document the distribution of employment across firm sizes, and estimate wage gaps that progressively condition on sector-year effects, worker characteristics, and formality. Because unobserved worker characteristics may be correlated with firm size, the estimates are conditional wage gaps rather than causal effects of firm size. This interpretation follows the broader firm-size wage premium literature: a premium that survives observable controls is informative about sorting, segmentation, firm-side pay differences, and possible allocative frictions, but it is not by itself proof of any one mechanism \citep{BrownMedoff1989,OiIdson1999,AbowdKramarzMargolis1999,CardHeiningKline2013,ElBadaouiStroblWalsh2010,BalkanTumen2016}.

The evidence delivers five main results. First, employment is concentrated at the bottom of the firm-size distribution: solo workers and firms with up to ten workers account for about two thirds of employment. Second, raw wage gaps are large across firm-size categories, sectors, education groups, and formality status, while the wage gap by sex is small by comparison. Third, the firm-size wage premium attenuates substantially after adding worker controls, sector-year fixed effects, and formality, but it does not disappear: in the full specification, workers in firms with 101 or more employees earn about 43 percent more than otherwise comparable solo workers. Fourth, the conditional premium is positive throughout the period; the point estimates suggest an increase between 2008 and 2025, but the change is not statistically significant. Fifth, formality status affects the firm-size wage premium. Formal workers earn more than informal workers within every firm-size category, but the conditional firm-size wage premium is much steeper among informal workers. Within formal employment, the conditional premium relative to formal solo workers is small, although workers in 101+ firms earn more than workers in intermediate-size formal firms.

The rest of the paper is organized as follows. Section~2 describes the GEIH data, the harmonization of firm-size categories, and the construction of real hourly labor income. Section~3 presents the descriptive evidence on employment shares and wage gaps by firm size, formality, education, gender, sector, and year. Section~4 introduces the econometric specifications. Section~5 reports the main regression results, the dynamic estimates, the formality interactions, and the elasticity benchmark, contrasting each result with the relevant literature. Section~6 discusses interpretation, limitations, and policy implications. Section~7 concludes.
